\section{Cài đặt Hệ thống}

\subsection{Backend}
Backend được xây dựng bằng FastAPI, đóng vai trò trung tâm trong 
kiến trúc hệ thống, chịu trách nhiệm xử lý các yêu cầu từ phía 
người dùng, quản lý nghiệp vụ, xác thực, phân quyền và kết nối 
với các dịch vụ như cơ sở dữ liệu, AI Service và hệ thống 
thanh toán.

\textbf{Cài đặt:}
\begin{itemize}[noitemsep]
    \item Khởi tạo môi trường ảo Python 3.11 và cài đặt các thư 
    viện cần thiết từ file \texttt{requirements.txt}.
    \item Cấu hình các biến môi trường trong file \texttt{.env} 
    bao gồm thông tin kết nối Neon PostgreSQL, Firebase, 
    Cloudinary và Stripe.
    \item Khởi động ngrok để tạo đường hầm công khai phục vụ 
    callback thanh toán và kết nối frontend.
    \item Khởi động server bằng script \texttt{start\_dev.ps1}, 
    hệ thống sẽ tự động cập nhật URL ngrok vào \texttt{.env} 
    và khởi chạy Uvicorn tại địa chỉ \texttt{localhost} cổng 8000.
    \item Kiểm tra các API tại \texttt{http://localhost:8000/docs}.
\end{itemize}

\subsection{Frontend}
Frontend được xây dựng bằng Nuxt.js 3 theo chế độ Single Page 
Application (SPA), đóng vai trò hiển thị giao diện người dùng, 
tiếp nhận tương tác và giao tiếp với backend thông qua các API.

\textbf{Cài đặt:}
\begin{itemize}[noitemsep]
    \item Cài đặt các thư viện cần thiết bằng lệnh 
    \texttt{yarn install}.
    \item Cấu hình kết nối API đến backend và các biến môi 
    trường Firebase trong file \texttt{.env}.
    \item Khởi động server phát triển bằng lệnh \texttt{yarn dev} 
    và truy cập hệ thống tại \texttt{http://localhost:3000}.
\end{itemize}

\subsection{Cơ sở dữ liệu}
Hệ thống sử dụng hai cơ sở dữ liệu song song phục vụ các mục 
đích khác nhau.

\textbf{Neon PostgreSQL} được sử dụng làm cơ sở dữ liệu quan hệ 
chính, lưu trữ thông tin người dùng, số đo cơ thể, garment, 
lịch sử thử đồ ảo, phiên chatbot và vector embeddings sản phẩm. 
Extension \textbf{pgvector} được kích hoạt để hỗ trợ tìm kiếm 
ngữ nghĩa phục vụ chatbot RAG.

\textbf{Firebase Firestore} được sử dụng để lưu trữ dữ liệu 
phi cấu trúc bao gồm sản phẩm, danh mục, đơn hàng, giỏ hàng, 
voucher và cấu hình giới thiệu cửa hàng. \textbf{Firebase 
Authentication} đảm nhận xác thực người dùng với các phương thức 
đăng nhập qua email và Google.

\textbf{DBeaver CE} được sử dụng làm công cụ quản lý và trực 
quan hóa dữ liệu trong Neon PostgreSQL, hỗ trợ truy vấn, 
chỉnh sửa và theo dõi dữ liệu trong quá trình phát triển 
\cite{dbeaver_docs}.

\subsection{Lưu trữ đám mây}
\textbf{Cloudinary} được sử dụng làm dịch vụ lưu trữ đám mây 
cho các tệp đa phương tiện bao gồm ảnh sản phẩm, ảnh người dùng 
và ảnh kết quả AI Virtual Try-On \cite{cloudinary_docs}. 
Cloudinary cung cấp khả năng quản lý, tối ưu hóa hình ảnh 
tự động và phân phối nội dung qua CDN, giúp cải thiện tốc độ 
tải trang.

\subsection{AI Service}
\textbf{FitDiT} được triển khai như một service độc lập, tích 
hợp vào backend thông qua module \texttt{fitdit.py}. Model được 
load lên GPU (NVIDIA RTX 4060 Laptop, 8GB VRAM) khi khởi động 
server và xử lý các yêu cầu thử đồ ảo qua endpoint 
\texttt{POST /tryon}.

\textbf{Ollama} được cài đặt trên máy cục bộ để chạy mô hình 
ngôn ngữ lớn \textbf{Qwen2.5 7B} phục vụ chatbot tư vấn thời 
trang. Chatbot kết hợp kỹ thuật RAG với pgvector để truy xuất 
thông tin sản phẩm phù hợp trước khi sinh phản hồi.

\textbf{Cấu hình Ollama:}
\begin{itemize}[noitemsep]
    \item Model: \texttt{qwen2.5:7b}
    \item Triển khai: cục bộ, không cần kết nối Internet
    \item Tích hợp: qua thư viện \texttt{ollama} 0.6 trong 
    FastAPI service
\end{itemize}

\subsection{Thanh toán}
Hệ thống tích hợp cổng thanh toán \textbf{Stripe} để xử lý 
thanh toán trực tuyến \cite{stripe_docs}.

\textbf{Cài đặt:}
\begin{itemize}[noitemsep]
    \item Cấu hình Stripe API key trong file \texttt{.env}.
    \item Sử dụng ngrok để công khai backend ra môi trường 
    Internet phục vụ Stripe webhook callback.
    \item Kiểm thử thanh toán trong môi trường sandbox trước 
    khi triển khai chính thức.
\end{itemize}

\subsection{Quy trình khởi động hệ thống}
\begin{enumerate}[noitemsep]
    \item Khởi động Ollama: \texttt{ollama serve} và load model 
    \texttt{qwen2.5:7b}.
    \item Khởi động Backend: chạy \texttt{.\textbackslash 
    start\_dev.ps1} tại thư mục 
    \texttt{D:\textbackslash Backend\_AR\_Clothes}.
    \item Chờ FitDiT load model lên GPU (khoảng 1--2 phút) 
    đến khi xuất hiện thông báo 
    \texttt{Application startup complete}.
    \item Khởi động Frontend: chạy \texttt{yarn dev} tại thư 
    mục \texttt{D:\textbackslash Web\_Clothes}.
    \item Truy cập hệ thống tại 
    \texttt{http://localhost:3000}.
\end{enumerate}
